% ============================================================
%  IACV 2025/26 Project — Synthetic Pose Estimation Pipeline
%  Mathematical Foundations of the Pipeline (steps 1-5)
%  Ready to compile on Overleaf (pdfLaTeX).
% ============================================================
\documentclass[11pt]{article}

\usepackage[utf8]{inputenc}
\usepackage[T1]{fontenc}
\usepackage{amsmath, amssymb, amsthm, mathtools}
\usepackage{bm}
\usepackage[margin=2.4cm]{geometry}
\usepackage{xcolor}
\usepackage{hyperref}
\usepackage{enumitem}

\hypersetup{colorlinks=true, linkcolor=blue!60!black, urlcolor=blue!60!black}

% --- Notation shortcuts ---
\newcommand{\R}{\mathbb{R}}
\newcommand{\SO}[1]{\mathrm{SO}(#1)}
\newcommand{\SE}[1]{\mathrm{SE}(#1)}
\newcommand{\Rgt}{\mathbf{R}_{\text{gt}}}
\newcommand{\tgt}{\mathbf{t}_{\text{gt}}}
\newcommand{\Pcam}{\mathbf{P}_{\text{cam}}}
\newcommand{\Pmark}{\mathbf{P}_{\text{mk}}}
\newcommand{\KK}{\mathbf{K}}

\title{Synthetic Pose Estimation Pipeline \\[0.3em]
       \large Mathematical Foundations --- Steps 1 to 5}
\author{IACV 2025/26 Project}
\date{\today}

\begin{document}
\maketitle
\tableofcontents
\vspace{1em}

% ============================================================
\paragraph{Overview}
% ============================================================

We synthesize the image a calibrated pinhole camera would observe of a
known planar square marker placed at a known pose.( no detection )

\begin{enumerate}[label=\textbf{Step \arabic*.}, leftmargin=*]
  \item Generate a 3D plane (the marker plane).
  \item Define a square on that plane of side $L$.
  \item Define a camera with known intrinsics $\KK$.
  \item Choose a known marker pose $(\Rgt, \tgt)$ relative to the camera.
  \item Project the four 3D corners into the image.
\end{enumerate}

Steps 1--5 are pure forward simulation (no estimation, no noise).
Steps 6-7 and the 2 experiments -> TO BE DONE

\subsection*{Notation}
\begin{itemize}[leftmargin=*]
  \item Vectors in bold lowercase: $\mathbf{t}$. Matrices in bold uppercase: $\mathbf{R}, \KK$.
  \item Points carry a subscript naming their frame:
        $\Pmark$ is in the marker frame, $\Pcam$ in the camera frame.
  \item Pixel coordinates: $\mathbf{p} = (u, v)^\top$.
  \item Coordinate frames are right-handed: in the camera frame, $X$ points right,
        $Y$ points down, $Z$ points into the scene.
\end{itemize}

% ============================================================
\section{Step 1 --- The 3D Marker Plane}
% ============================================================

The marker is a flat object, so the four corners are coplanar. We attach a
coordinate frame to the marker (the \emph{marker frame}) with the plane
chosen as $Z = 0$:
\begin{equation}
  \text{marker plane} \;=\; \{\, (X, Y, 0) \,:\, X, Y \in \R \,\}.
\end{equation}
The origin is placed at the geometric center of the future square. The
$X$ axis points right, $Y$ points up (in the plane), and $Z$ is the outward
normal to the marker surface.

\paragraph{Why $Z = 0$ matters.}
A general projection of a 3D point requires the full $3\times 3$ rotation
matrix. If the third coordinate is identically zero, the third column of
$\mathbf{R}$ multiplies $0$ and drops out:
\begin{equation}
  \mathbf{R} \begin{pmatrix} X \\ Y \\ 0 \end{pmatrix}
  \;=\; X \mathbf{r}_1 + Y \mathbf{r}_2,
\end{equation}
where $\mathbf{r}_1, \mathbf{r}_2$ are the first two columns of $\mathbf{R}$. This
collapse is the structural reason a single $3\times 3$ matrix --- a
\emph{homography} --- captures all the geometry of a planar marker, which
is what later pose-estimation steps will exploit.

% ============================================================
\section{Step 2 --- The Square on the Plane}
% ============================================================

On the plane $Z = 0$ we place a square of side $L$, centered at the marker
origin and axis-aligned to the marker frame. The four corner coordinates,
in the order used by OpenCV's \verb|SOLVEPNP_IPPE_SQUARE|, are
\begin{equation}
  \Pmark^{(0)} = \begin{pmatrix} -L/2 \\ +L/2 \\ 0 \end{pmatrix}, \;
  \Pmark^{(1)} = \begin{pmatrix} +L/2 \\ +L/2 \\ 0 \end{pmatrix}, \;
  \Pmark^{(2)} = \begin{pmatrix} +L/2 \\ -L/2 \\ 0 \end{pmatrix}, \;
  \Pmark^{(3)} = \begin{pmatrix} -L/2 \\ -L/2 \\ 0 \end{pmatrix}.
\end{equation}
The parenthesized superscript is the corner index; the subscript ``mk''
indicates the marker frame.

\paragraph{Corner ordering matters.}
The order of the 3D list must match the order of the 2D observations,
otherwise the correspondences are scrambled and the algorithm finds the
wrong pose (or fails). The order above (counter-clockwise starting at
top-left, viewed from $+Z$) is the convention shared by OpenCV's ArUco
detector and \verb|SOLVEPNP_IPPE_SQUARE|.

% ============================================================
\section{Step 3 --- The Pinhole Camera}
% ============================================================

\subsection{Geometric model}

A pinhole camera maps points in the camera frame to points on a 2D image.
The camera origin is at the world origin; the image plane sits at $Z = f$,
where $f$ is the focal length. By similar triangles, a 3D point
$\Pcam = (X, Y, Z)^\top$ projects to
\begin{equation}
  x = f\cdot\frac{X}{Z}, \qquad y = f\cdot\frac{Y}{Z}.
  \label{eq:similar-triangles}
\end{equation}
The division by $Z$ is the \emph{perspective divide}: it is the only
non-linearity in the model and is responsible for foreshortening (farther
objects look smaller).

\subsection{The intrinsic matrix \texorpdfstring{$\KK$}{K}}

Real images are measured in pixels. The focal length is encoded as
$f_x, f_y$ (in pixels, possibly different along the two axes if the pixel
elements are not square), and the optical axis hits the image at the
\emph{principal point} $(c_x, c_y)$. The pixel coordinates are
\begin{equation}
  u = f_x \cdot \frac{X}{Z} + c_x, \qquad
  v = f_y \cdot \frac{Y}{Z} + c_y.
  \label{eq:pixel-projection}
\end{equation}
Collected into a matrix:
\begin{equation}
  \boxed{\quad
    \KK \;=\; \begin{pmatrix} f_x & 0 & c_x \\ 0 & f_y & c_y \\ 0 & 0 & 1 \end{pmatrix}
  \quad}
\end{equation}
so that the projection has the compact homogeneous form
\begin{equation}
  \tilde{\mathbf{p}} \;=\; \KK\,\Pcam,
  \qquad
  \mathbf{p} \;=\; \begin{pmatrix} \tilde p_1/\tilde p_3 \\ \tilde p_2/\tilde p_3 \end{pmatrix}.
\end{equation}
The first is a linear matrix product, the second is the perspective divide.

\subsection{Normalized image coordinates}

Apply $\KK^{-1}$ to a pixel $\mathbf{p}$:
\begin{equation}
  \begin{pmatrix} x_n \\ y_n \\ 1 \end{pmatrix}
  \;=\; \KK^{-1}
  \begin{pmatrix} u \\ v \\ 1 \end{pmatrix}
  \;=\;
  \begin{pmatrix} (u - c_x)/f_x \\ (v - c_y)/f_y \\ 1 \end{pmatrix}.
\end{equation}
The pair $(x_n, y_n)$ are the \emph{normalized image coordinates}: the
location where the ray through pixel $\mathbf{p}$ crosses the plane
$Z = 1$ in camera space. They are dimensionless (no pixels, no focal
length) and encode pure ray directions. Pose-estimation algorithms
work in this representation so that the camera-specific quantities $\KK$
do not contaminate the geometric reasoning.

% ============================================================
\section{Step 4 --- Ground-Truth Pose}
% ============================================================

A point in the marker frame is brought into the camera frame by a
rigid-body transform:
\begin{equation}
  \boxed{\quad \Pcam \;=\; \mathbf{R}\,\Pmark + \mathbf{t} \quad}
  \label{eq:rigid-transform}
\end{equation}
with $\mathbf{R} \in \SO{3}$ and $\mathbf{t} \in \R^3$. The pair
$(\mathbf{R}, \mathbf{t})$ is an element of the Special Euclidean group
$\SE{3}$, and is what we mean by \emph{pose}.

\subsection{The rotation group \texorpdfstring{$\SO{3}$}{SO(3)}}

A $3\times 3$ matrix $\mathbf{R}$ is a rotation if and only if
\begin{equation}
  \mathbf{R}^\top \mathbf{R} \;=\; \mathbf{I},
  \qquad
  \det(\mathbf{R}) \;=\; +1.
  \label{eq:so3-constraints}
\end{equation}
The first condition (orthonormal columns) preserves lengths and angles; the
second excludes reflections, which would flip a right-handed frame into a
left-handed one.

\paragraph{Degrees of freedom.}
A $3\times 3$ matrix has $9$ entries. The orthogonality condition
$\mathbf{R}^\top \mathbf{R} = \mathbf{I}$ imposes $6$ independent constraints. That
leaves \textbf{3 DOF} --- the minimum number of parameters needed to
describe a rotation in 3D.

\subsection{Representation 1: Euler angles}

Decompose any rotation as three elementary rotations around the coordinate
axes. The \emph{intrinsic xyz} convention used in our code is
\begin{equation}
  \mathbf{R}(\alpha,\beta,\gamma) \;=\; \mathbf{R}_x(\alpha)\,\mathbf{R}_y(\beta)\,\mathbf{R}_z(\gamma),
\end{equation}
where each $\mathbf{R}_\bullet(\cdot)$ is a $3\times 3$ planar rotation, e.g.
\begin{equation}
  \mathbf{R}_x(\alpha) =
  \begin{pmatrix} 1 & 0 & 0 \\ 0 & \cos\alpha & -\sin\alpha \\ 0 & \sin\alpha & \cos\alpha \end{pmatrix}.
\end{equation}
\textbf{Pros:} human-intuitive --- $(15^\circ, 10^\circ, 5^\circ)$ is easy to picture.\\
\textbf{Cons:} suffer from \emph{gimbal lock} (a configuration where two
axes align and one DOF is lost) and from non-unique decompositions. We use
Euler angles only to \emph{specify} the ground-truth pose conveniently; we
never use them for optimization.

\subsection{Our concrete ground-truth pose}

In the code we choose
\begin{align}
  \text{Euler (xyz, deg)} & : \quad (\alpha, \beta, \gamma) = (15,\,10,\,5) \\
  \tgt & : \quad (0.05,\, 0.02,\, 0.50)^\top \;\text{m}
\end{align}
which yields
\[
  \Rgt \approx
  \begin{pmatrix}
    +0.9811 & -0.0394 & +0.1897 \\
    +0.0858 & +0.9662 & -0.2432 \\
    -0.1736 & +0.2549 & +0.9513
  \end{pmatrix}.
\]
Note that the total rotation angle ($\approx 18.33^\circ$) is
\emph{not} $15^\circ + 10^\circ + 5^\circ$ because Euler rotations do not
compose by simple addition.

% ============================================================
\section{Step 5 --- Forward Projection}
% ============================================================

Steps 1--4 fix all the inputs. Step 5 chains them into the observed image
points. For each marker corner $\Pmark^{(i)}$:
\begin{equation}
  \boxed{\quad
    \mathbf{p}^{(i)}
    \;=\;
    \pi\!\Big(\, \KK\,(\,\Rgt \Pmark^{(i)} + \tgt\,) \,\Big)
  \quad}
  \label{eq:full-projection}
\end{equation}
where the projection $\pi:\R^3 \to \R^2$ is the perspective divide
\begin{equation}
  \pi\!\begin{pmatrix} \tilde x \\ \tilde y \\ \tilde z \end{pmatrix}
  \;=\;
  \begin{pmatrix} \tilde x / \tilde z \\ \tilde y / \tilde z \end{pmatrix}.
\end{equation}

\paragraph{In the code, Step 5 is split in two sub-steps:}
\begin{enumerate}[label=(\alph*)]
  \item \textbf{Rigid transform.} $\Pcam^{(i)} \;=\; \Rgt \Pmark^{(i)} + \tgt$.
        Brings each marker corner from the marker frame into the camera frame.
  \item \textbf{Pinhole projection.} $\mathbf{p}^{(i)} \;=\; \pi(\KK\,\Pcam^{(i)})$.
        Maps the camera-frame point to pixel coordinates via
        \eqref{eq:pixel-projection}.
\end{enumerate}

\paragraph{Possible sanity checks at this point:}
\begin{itemize}[leftmargin=*]
  \item All $\Pcam^{(i)}$ have $Z > 0$ (in front of the camera).
  \item All $\mathbf{p}^{(i)}$ fall inside $[0, W) \times [0, H)$ (visible in the image, W as width, H as height).
  \item The four $\mathbf{p}^{(i)}$ form a \emph{convex quadrilateral} that looks
        like a perspective-warped square.
\end{itemize}

\end{document}
